% Part: applied-modal-logics
% Chapter: temporal-logic
% Section: introduction

\documentclass[../../../include/open-logic-section]{subfiles}

\begin{document}

\olfileid{aml}{tl}{int}

\olsection{سريزه}

زماني منطقونه د هغو دعوو څېړنه کوي چې يو څۀ به په راتلونکي کښې واقع شي يا
په تېر وخت کښې واقع شوي وي۔ د زماني منطق بنسټ ايښودونکے آرتر پرايور ګڼل
کېږي؛ هغه ورته \emph{د زمانو منطق} وايۀ۔ دلته به د زماني منطق په څېړنه کښې
تر ډېره د پرايور هغه لومړنۍ مودالي طريقه تعقيبوو چې د قضيېيز منطق بنسټيز
چوکاټ ته زماني عاملونه ورزياتوي۔ قضيېيز منطق عموماً دعوې داسې ګڼي چې زمان
ئې نۀ وي ټاکل شوے۔

د بېلګې په توګه، زه ښايي د بيزي (Beezie) په نوم د يوه سپي په اړه خبرې
وکړم چې کله کښېني او کله نۀ کښېني، لکه د سپيو عادت چې وي۔ په کلاسيک
منطق کښې دا ويل متناقض دي چې بيزي ناست دے او بيزي ناست نۀ دے۔ خو ښکاره
ده چې دواړه خبرې رښتيا کېدے شي، البته په يوه وخت کښې نۀ۔ ژبې ته د زماني
عاملونو په ورزياتولو سره دا خبره په نسبتاً اسانه توګه څرګندولے شو۔ د
زماني عاملونو ورزياتول موږ ته دا هم راښيي چې ولې له «بيزي به انعام يا
توپ ترلاسه کړي» څخه «بيزي به انعام ترلاسه کړي يا بيزي به توپ ترلاسه
کړي» ته استدلال معتبر دے۔

خو زماني منطق ډېرې فلسفي پوښتنې راولاړوي چې ښايي موږ د زماني منطق يوه
چوکاټ پر بل غوره کولو ته وهڅوي۔ د بېلګې په توګه، راتلونکے اتفاقي
بيان د راتلونکي په اړه هغه بيان دے چې نۀ لازم دے او نۀ ناممکن۔ که ووايو
«ريچارډ (Richard) به سبا د خوراکي توکو دوکان ته لاړ شي»، نو د يوه داسې
کار په اړه دعوه کوو چې لا نۀ دے شوے او د رښتيا ارزښت ئې د بحث وړ دے۔
ان دا هم د بحث وړ ده چې آيا په لومړي سر کښې دې دعوې ته د رښتيا ارزښت
\emph{ټاکل} کېدے شي که نۀ۔ که موږ سخت جبرپال يو، نو ښايي دا ومنو چې
دا جمله په حقيقت کښې رښتيا يا دروغ ده، ان تر دې چې ياد کار وشي مخکښې؛
يوازې ښايي موږ لا د هغې د رښتيا ارزښت نۀ پېژنو۔ برعکس، ښايي موږ پر
رښتيا پرانيستي راتلونکي باور ولرو چې پکښې د راتلونکو اتفاقي بيانونو
د رښتيا ارزښتونه لا ناتعيين وي۔

په پايله کښې، د وخت د جوړښت او ماهيت په اړه ډېرې داسې ژمنې د زماني
منطقونو د مدلونو او چوکاټونو په ټاکلو کښې نغښتې وي۔ د بېلګې په توګه،
له ځانه پوښتلے شو چې آيا داسې مدلونه جوړ کړو چې وخت پکښې خطي، څوڅانګيز
يا ان دوراني وي۔ ښايي دا هم وټاکو چې زموږ زماني مدلونه د پيل او پاې
ټکي لري که نۀ، او وخت د بېلو شېبو په توګه وښيو که د يوه پيوسته امتداد
په توګه۔

\end{document}
